\documentclass[11pt]{article}

\usepackage[letterpaper,margin=1in]{geometry}
\usepackage{microtype}
\usepackage{amsmath}
\usepackage{amssymb}
\usepackage{booktabs}
\usepackage{array}
\usepackage{enumitem}
\usepackage{float}
\usepackage[dvipsnames]{xcolor}
\usepackage[most]{tcolorbox}
\usepackage[hidelinks]{hyperref}
\usepackage{dsfont}

\emergencystretch=1.5em

\definecolor{boxblue}{HTML}{2C3E50}
\newtcolorbox{keybox}{enhanced,breakable,colback=boxblue!5,colframe=boxblue!55,
  boxrule=0.5pt,arc=2pt,left=8pt,right=8pt,top=5pt,bottom=5pt}

\newtcolorbox{plain}{
  enhanced, breakable,
  colback=Sepia!7, colframe=Sepia!55, boxrule=0.4pt, arc=2pt,
  left=7pt, right=7pt, top=4pt, bottom=4pt,
  fontupper=\small,
  coltitle=Sepia!72!black, fonttitle=\bfseries\sffamily\footnotesize,
  title={IN PLAIN TERMS}
}

\newcommand{\code}[1]{\texttt{#1}}
\newcommand{\rfckw}[1]{\textsc{#1}}
\newcommand{\ind}{\mathds{1}}

\title{\bfseries A Scalable Trigger-Metadata Model for\\
Internet Reachability under the FedRAMP VDR/VER Rules}
\author{Matthew Venne \\ \small Chief Technology Officer, stackArmor}
\date{July 2026}

\begin{document}
\maketitle

\begin{abstract}
\noindent
The FedRAMP Vulnerability Detection and Response (VDR) and Vulnerability
Evaluation and Reporting (VER) rules require providers to determine, per
finding, whether a vulnerability is internet-reachable. Direct exposure is
often computable from information providers already possess: public routes,
load-balancer listeners, firewall policy, running processes, affected
software, and CVSS attack vector. Transitive reachability is different.
Vulnerability sources and scanners rarely publish structured metadata
describing whether a CVE can be triggered by internet-originating content after
that content crosses a component boundary. Without that upstream fact,
providers face two bad defaults: classify every vulnerability on every
internet-touched component as reachable, destroying prioritization, or
approximate reachability with direct accessibility and risk missing exceptional
indirect triggers such as Log4Shell. Deployment-specific data-flow evidence can
refine the second step---which affected components actually receive the content
and what happens on the way---but vulnerability tracking systems do not
generally integrate that evidence either.

This paper recommends an operational model between those extremes. It computes
direct finding-level reachability and adds a
reusable CVE-level \code{indirect\_internet\_trigger} promotion when vendor
guidance, threat intelligence, exploit analysis, agentic enrichment, or
analyst review establishes a transitive trigger. A coarse asset assertion
identifies components that receive internet-derived content. Existing
assessment and continuous-monitoring evidence supplies a standing, revocable
prevention assertion for the input-control classes it actually tests. For
Rev5 services, that evidence already includes recurring web-application
scanning, remediation, and assessor validation; it does not require a new
application-wide data-flow graph. This is the paper's present recommendation,
not an interim step contingent on a future provider capability. It causes the
operational IRV set to converge toward the directly accessible surface, plus
known indirect triggers not covered by the standing assertion. A separate
full-information equation documents how
structured trigger profiles and optional deployment-edge evidence could refine
the result if the vulnerability-intelligence ecosystem later supplies them. It
is an interoperability reference, not a proposed FedRAMP evidence requirement.
Agentic enrichment is treated as a controlled bridge over missing upstream
metadata, not as an authoritative replacement for it.
\end{abstract}

\section{Introduction}
\label{sec:intro}
FedRAMP's 2026 rules require providers to evaluate detected vulnerabilities in
the context of the cloud service offering and report whether each is an
Internet-Reachable Vulnerability (IRV), whether it is a Likely Exploitable
Vulnerability (LEV), and the Potential Agency Impact N-rating (PAIN) of
exploitation. Those answers select the response timeframe under
\code{VDR-TFR-PVR}; the IRV result is also included in the vulnerability detail
report required by \code{VER-RPT-VDT}. The rules deliberately use
\emph{internet-reachable}, not merely \emph{internet-accessible}: a vulnerable
resource may sit on a private network and still process a payload that
originated on the public internet.\cite{fedramp-definitions,fedramp-ver}

That distinction is technically correct and operationally difficult. A
scanner can usually identify a CVE and affected component. It may supply a
CVSS vector, a CWE, a short description, and references. It generally does not
state, in machine-readable form, whether attacker-controlled content can cross
a process boundary and still trigger the vulnerability, which content format
or processing operation is required, or whether the attacker must remain the
connection peer. Nor does the vulnerability record normally know the
provider's application graph: which API writes which queue, which worker parses
which upload, or which validation is applied before a database query.

The missing information cannot be recreated by multiplying the number of
asset-level scanner findings. If one CVE is present on ten thousand assets,
the trigger mechanism is still chiefly a fact about the CVE and affected
product family. It should be determined once, carried with provenance, and
reused. Deployment applicability should likewise be evaluated by component
and path archetype rather than by repeating the same analysis for every
instance. FedRAMP expressly permits logical grouping of affected resources to
improve evaluation and response efficiency under \code{VER-EVA-GRV}.\cite{fedramp-ver}

This paper defines a method that can be implemented now and separately records
optional extension points for upstream standardization. Sections~\ref{sec:two-defaults}
and~\ref{sec:gaps} define the operational problem. Section~\ref{sec:model}
gives the recommended equation first and the full-information reference form
second. Sections~\ref{sec:direct} through~\ref{sec:evidence-plane} describe the
inputs. The remaining sections apply the method to triage, controls, worked
examples, and validation.

\begin{plain}
FedRAMP asks a question scanners cannot fully answer from their normal output.
Scanners are reasonably good at saying what software is present and which
network service is exposed. They are not designed to say whether a string,
file, message, or database record can carry a particular exploit through three
application tiers. The solution is not to perform the same manual research for
every asset. It is to enrich the vulnerability once, combine that result with
reusable deployment and standing control evidence, and preserve uncertainty
when either side is missing.
\end{plain}

\section{The Metadata Gap Creates Two Bad Defaults}
\label{sec:two-defaults}
FedRAMP defines an IRV as a vulnerability that might be exploited or triggered
by a payload originating on the public internet. Its notes expressly include
resources with no direct internet route when they receive a payload or take an
action triggered by internet activity. The evaluation guidance also says that
payload interception, filtering, sanitization, rejection, or comparable
prevention before vulnerable processing can remove IRV status.\cite{fedramp-definitions,fedramp-ver}

Those rules establish the right semantic boundary: the unit is the
\emph{finding}, and the question is whether a qualifying trigger can reach the
vulnerable processing operation. They do not, however, provide the metadata
needed to automate that test. When a provider sees a local parser CVE on a
queue worker, neither the scanner finding nor the CVSS vector necessarily says
whether a public file upload can activate that parser. FIRST's CVSS guidance
even notes that malicious network data passed from one system to a separate
vulnerable system can be scored \code{AV:L}; base attack vector alone cannot
resolve the transitive question.\cite{cvss-user-guide}

The gap creates two operational extremes.

\paragraph{Maximal propagation.}
Assume every vulnerability on every component that ever receives
internet-derived data might be triggered by that data. This preserves recall,
but a connected cloud service quickly approaches universal IRV classification:
requests reach applications, records reach databases, messages reach workers,
and strings reach logs. The IRV/NIRV distinction then contributes little
prioritization while the provider still has to record and report the result.

\paragraph{Accessibility-only approximation.}
Evaluate only direct network exposure. This is scalable and reproducible, but
it misses known transitive shapes such as a crafted string reaching a logging
library or a malicious file reaching an interior parser. It treats the absence
of metadata as evidence that no indirect trigger exists, which is not the same
claim.

The practical answer today is an explicitly metadata-constrained
approximation: compute direct reachability from available evidence, promote
known indirect-trigger CVEs through a reusable enrichment, and apply the
standing prevention evidence already produced through assessment and
continuous monitoring. In set terms,

\[
R_{\mathrm{operational}}=D\cup(K\setminus\mathcal P_0),
\]

where $D$ is the set of directly triggerable findings and $K$ is the set of
findings promoted by known indirect-trigger intelligence. $\mathcal P_0$ is the
subset for which $P_0(v,a)=1$ under a current assessment- and ConMon-backed
prevention assertion. The
true set may also contain triggerable members of an unresolved population $U$
and depends on whether the asserted controls actually prevent the trigger:

\[
R_{\mathrm{true}}=D\cup(K\setminus\mathcal P_{\mathrm{true}})
\cup U_{\mathrm{triggerable}}.
\]

Under current metadata constraints, $R_{\mathrm{operational}}$ is therefore
dominated by the directly accessible surface plus known, uncovered promotions.
This is the paper's convergence claim: not that reachability and accessibility
are semantically equivalent, but that recurring assessment, DAST, and
remediation support a rebuttable operational presumption that tested downstream
trigger classes are prevented. A failed, stale, out-of-scope, or contradicted
assertion does not enter $\mathcal P_0$.

\begin{plain}
There are two easy answers, and both are bad. Calling every flaw touched by
internet data reachable makes the label almost useless. Looking only at open
internet services misses the rare but important flaws whose trigger rides
inside data. Today's workable compromise starts with the direct surface and
adds known indirect cases once per CVE. Existing assessment and recurring web
testing then remove the cases for which the relevant input control is currently
supported. This makes the answer look much like accessibility in steady state,
with known uncovered triggers as explicit exceptions. The missing cases remain
a disclosed limitation, and better upstream metadata is the durable fix.
\end{plain}

\section{Current Inputs and Optional Refinement}
\label{sec:gaps}
The recommended method joins two kinds of information at the resolution
available today. Trigger status is chiefly a vulnerability-level fact;
deployment applicability is a provider-level fact. Neither requires a provider
to construct a complete application-wide taint or data-flow proof.

\subsection{Vulnerability trigger metadata}
The vulnerability side must describe how the flaw is activated, not merely
where the affected package is installed. Useful fields include the carrier of
the attacker input, its format and encoding, the operation performed on it,
the degree of attacker influence required, and whether the attacker must
remain the connection peer. This information should originate with vendors,
CVE Numbering Authorities, coordinated-disclosure participants, or other
authoritative vulnerability publishers, then be distributed through scanners
and threat-intelligence feeds.

Today's scanners are a natural delivery point but not necessarily the
authoritative source. Their ordinary findings often contain only a compact
description whose purpose is detection and remediation, not a complete exploit
contract. A CWE can route research toward a family of mechanisms; it cannot
stand in for the exact trigger of a CVE.

\subsection{Deployment applicability and prevention evidence}
The current profile needs governed assertions about which component archetypes
receive internet-derived content and which tested trigger classes are prevented
before vulnerable processing. FedRAMP assessment and ConMon already produce a
recurring evidence stream for the latter: web-application scan scope and
configuration, results, remediation and retest history, control implementation
evidence, and assessor validation. A provider may refine those assertions with
evidence about how content moves through the service and what
each boundary does to it. A higher-resolution representation can include
API-to-service calls, queue publication and consumption, database and object
storage flows, logging paths, file-processing pipelines, protocol translation,
and validation or sanitization behavior. Each claim should be bound to the
source revision, built artifact, deployed configuration, and evidence that
supports it.

Individual ingredients exist today: infrastructure as code, service catalogs,
flow telemetry, API specifications, SAST results, unit tests, integration
tests, and CI/CD provenance. But the industry lacks a broadly adopted,
machine-readable schema and assurance method for combining those artifacts
into an independently verifiable, vulnerability-specific reachability record
that a vulnerability-management platform can ingest, map to affected findings,
and re-evaluate when the supporting code, configuration, topology, or evidence
changes. This paper treats such a record as an optional interoperability
refinement, not as a prerequisite for using the recommended method.

\begin{keybox}
\centering
\textbf{Vulnerability Trigger Profile}
\quad$+$\quad
\textbf{Deployment and Standing Control Evidence}
\quad$=$\quad
\textbf{Finding-Level Reachability}
\end{keybox}

\begin{plain}
A provider can usually identify broad groups of components that receive
internet-derived data---APIs, queues, workers, databases, and logging systems.
FedRAMP assessment and recurring web scans also provide a standing assertion
that the input-control classes in their validated scope are working and that
findings are remediated. That is enough for a scoped triage presumption; it is
not proof against an unknown trigger or an untested path. Proving the exact
path, what happens at each handoff, and which deployed code implements a
validation control is harder. The pieces exist across
infrastructure code, telemetry, tests, and CI/CD records, but vulnerability
systems cannot yet combine them into a standard, continuously valid
reachability record. The recommended method therefore uses governed component
and control archetypes today. A higher-resolution claim still needs a verified
prevention point on every relevant path before the vulnerable operation; not
every edge must sanitize the data.
\end{plain}

\section{Two-Surface Finding-Level Model}
\label{sec:model}
The original network descriptor remains useful, but network conversations and
relayed content are different objects. The model therefore uses two linked
surfaces rather than forcing both into one tuple.

\subsection{The network surface}
Let $N(a)$ be the set of network descriptors directly deliverable to component
$a$, and let $X_N(v)$ be the network descriptors under which vulnerability $v$
can be directly triggered. A network descriptor is

\[
n=\langle\kappa,\tau,\pi,\nu,\delta\rangle,
\]

where $\kappa$ is the CVSS attack-vector class, $\tau$ is transport and port,
$\pi$ is application protocol, $\nu$ is protocol version or ALPN, and $\delta$
is direction. Direct reachability is

\begin{equation}
D(v,a)=\ind\!\left[N(a)\cap X_N(v)\neq\varnothing\right].
\label{eq:direct}
\end{equation}

CVSS \code{AV:N} is evidence that $X_N(v)$ may contain a network trigger.
\code{AV:L}, \code{AV:A}, and \code{AV:P} exclude a vulnerability from this
\emph{direct} branch; they do not establish that $v$ has no indirect
content-trigger path.

\subsection{The recommended operational profile}
Structured per-CVE content-trigger profiles and per-edge transformation graphs
are not generally available in current scanners and vulnerability tracking
systems. The method recommended by this paper uses three available, governed
predicates:

\begin{itemize}[leftmargin=1.5em,nosep]
\item $J(v)$ is the CVE-level
  \code{indirect\_}\allowbreak\code{internet\_}\allowbreak\code{trigger}
  status: confirmed, ruled out, or unknown.
\item $C_0(a)$ asserts that component $a$ receives internet-derived content,
  derived from an asset archetype, an imported or observed flow, or a governed
  operator assertion.
\item $P_0(v,a)$ asserts that current assessment and ConMon evidence covers the
  relevant trigger or control class and supports prevention before vulnerable
  processing. It is resolved by application or deployment archetype rather than
  independently for every asset.
\end{itemize}

The recommended equation is

\begin{keybox}
\begin{equation}
\mathrm{IRV}_0(v,a)=
D(v,a)
\;\lor\;
\left(
\ind[J(v)=\mathrm{confirmed}]\cdot C_0(a)
\cdot\left(1-P_0(v,a)\right)
\right).
\label{eq:operational}
\end{equation}
\end{keybox}

The status \code{unknown} is retained as data even though the baseline
classifier treats it as not promoted. This is a deliberate scalability
presumption, not proof of non-reachability. Vendor guidance, a CNA update,
threat intelligence, exploit analysis, agentic enrichment, or analyst review
can promote $J(v)$ once and apply the result across every affected finding.

$P_0(v,a)=1$ is a standing assurance claim, not a mathematical proof. It
requires current evidence of the relevant scan or test scope, assessor-validated
configuration, remediation and successful retest of relevant findings, and no
known bypass or material change that invalidates the evidence. A confirmed
trigger outside the tested class, a new relevant finding, a failed retest,
missing or stale coverage, or an alternate path sets $P_0(v,a)=0$ until the
record is re-established. This makes the presumption explicit and auditable
instead of silently assuming that all downstream content is safe.

Equation~\ref{eq:operational} is not a placeholder for a provider-side system
that FedRAMP expects but industry has not yet built. It is the operational
classifier this paper recommends under the current rules and metadata
environment. Its CVE-level promotion makes the exceptional transitive cases
explicit, while $P_0$ reuses evidence providers and assessors already handle.
Neither requires every provider to rediscover an unpublished trigger or prove
every application data-flow edge.

\subsection{Full-information reference extension (non-normative)}
For interoperability and upstream metadata design, it is useful to show
how more structured inputs fit without changing the recommended classifier's
logic. Let $C(a)$ be the set of internet-originating content capabilities that
reach component $a$, and let $X_C(v)$ be the content trigger profiles capable
of triggering $v$. A content descriptor is

\[
c=\langle\gamma,\phi,o,\iota,b\rangle,
\]

where $\gamma$ is the carrier (request field, uploaded file, queue message,
database record, or log entry), $\phi$ is format and encoding, $o$ is the
processing operation (parse, deserialize, query, log, render, execute, or
decompress), $\iota$ is the degree of attacker influence (data-only,
structural/control, or arbitrary bytes), and $b$ records whether exploitation
requires the attacker to remain the connection peer.

Each documented data-flow edge $e$ may apply a transformation
$T_e:2^{\mathcal C}\to2^{\mathcal C}$. For a path $p$ from a public entry point
to $a$,

\begin{equation}
C(a)=\bigcup_{p\in\mathrm{paths}(\mathrm{internet},a)}
\left(T_{e_n}\circ\cdots\circ T_{e_1}\right)(C_{\mathrm{internet}}).
\label{eq:content-surface}
\end{equation}

The full-information reference form is

\begin{keybox}
\begin{equation}
\mathrm{IRV}_{F}(v,a)=
\ind\!\left[
N(a)\cap X_N(v)\neq\varnothing
\;\lor\;
C(a)\cap X_C(v)\neq\varnothing
\right].
\label{eq:fullinfo}
\end{equation}
\end{keybox}

Equation~\ref{eq:fullinfo} is a reference semantics for publishers and tools
that elect to supply richer inputs. It is not a claim that current providers
possess those inputs, a recommendation that they build them now, or a new
evidence obligation inferred from FedRAMP. Better upstream data may refine
results over time without making Equation~\ref{eq:operational} unreasonable or
incomplete for its declared operational purpose.

\begin{plain}
There are two ways an exploit reaches code. It can arrive as a network
conversation, or it can ride inside content handed from one component to the
next. We can automate much of the first path today. For the second, today's
system uses a CVE-level promotion, a coarse ``receives internet content'' tag,
and the standing prevention evidence already maintained through assessment and
ConMon. The full-information form shows how optional structured metadata can
plug in later; it does not turn that optional refinement into today's
requirement.
\end{plain}

\section{Computing the Direct Network Branch}
\label{sec:direct}
The direct branch is the mature part of the model. It is an OR across every
public ingress path and an AND of conditions within each path:

\[
N(a)=\bigcup_{P\in\mathrm{paths}_{\mathrm{direct}}(a)}N_P(a).
\]

A path contributes descriptors only when a route reaches a live listener and
the relevant protocol/version survives the controls on that path. The
following gates are implementation procedures, not independent definitions of
IRV.

\subsection{Gate 1: Addressing and routing}
A direct path requires a public address or public ingress, an active route,
and forwarding to the component. A public host path normally requires a public
IP and an internet-gateway route. A load-balanced backend requires neither on
the backend itself; the public load balancer and its target binding establish
the path. Egress-only NAT does not create an inbound path.

\subsection{Gate 2: Source restrictions}
Firewall and load-balancer source restrictions determine who can attempt the
path and should be recorded. Attribution and prefix breadth are useful triage
signals, but they do not automatically prove NIRV: traffic from an attributed
partner range can still originate on the public internet, and a compromised
authorized source can still send a trigger. A source restriction removes the
descriptor only when the provider's documented interpretation and evidence
establish that the qualifying public source population cannot use the path.
Otherwise it informs LEV and residual likelihood.

An aggregate source range broader than a \code{/20} may remain a useful review
threshold, but it is a heuristic for examining allowlist quality, not a
reachability boundary.

\subsection{Gate 3: Edge authorization boundaries}
A separate, bypass-free remote-access boundary can remove a backend from the
direct internet surface. The method removes a path from $N(a)$ when an IAP,
VPN, or comparable edge component authorizes the connection before traffic
reaches the backend; fronts organizational personnel and tooling; is the
backend's sole ingress path; and prevents direct invocation or spoofing around
the boundary. The edge component remains directly internet-accessible, but the
protected backend is NIRV for that direct path. This is a path decision, not a
claim that the boundary sanitizes payload content; a separately confirmed
indirect content path is still evaluated under $J(v)$ and $C_0(a)$.

Application-level authentication is not an IRV/NIRV gate in this method. By
the time an application's login, session, or authorization logic runs, traffic
has reached the component. Application authentication may inform LEV or an
attested mitigation. It may support a lower PAIN rating only when evidence
shows that the enforced role, tenant scope, or privilege boundary reduces the
customer effect of successful exploitation; the existence of a login alone
does not lower PAIN.

\subsection{Gate 4: Kubernetes and container routing}
Container exposure is determined from the union of public Gateway API routes,
Ingress resources, public \code{LoadBalancer} services, reachable
\code{NodePort} services, and \code{hostNetwork}/\code{hostPort} bindings.
Cluster state must be joined with cloud routing and firewall state: Kubernetes
alone cannot say whether a node port is reachable from the internet. Missing
either evidence source leaves the path unresolved.

\subsection{Gate 5: Process, listener, protocol, and version}
An exposed host does not make every installed package directly reachable. The
component must map to a running process and a reachable listener that speaks a
protocol and version in $X_N(v)$. Package inventory joined with process and
socket ownership supplies the strongest direct mapping. Protocol termination
also matters: an HTTP/2 flaw on a backend is not directly reachable through a
load balancer that terminates HTTP/2 and speaks only HTTP/1.1 to that backend,
provided no alternate HTTP/2 path exists.

No-listener evidence closes only the direct branch. A socketless worker can
remain indirectly reachable if a confirmed trigger rides inside a file,
message, record, or log entry it processes.

\section{Proposed Vulnerability Trigger Profiles}
\label{sec:trigger-profile}
The Vulnerability Trigger Profile is a proposed missing ecosystem artifact, not
a record that current scanners, CNAs, or vulnerability-management platforms
consistently publish or consume. A provider can implement the coarse $J(v)$
status today as a custom field or enrichment, but the structured profile below
has no broadly adopted interchange or assurance standard. Its purpose is to
show what upstream publishers and scanners would need to supply so the indirect
branch can be automated consistently. A profile may contain several trigger
alternatives because a CVE can have different mechanisms under different
configurations.

\begin{verbatim}
cve: CVE-202x-xxxxx
product: example-parser
indirect_trigger:
  status: confirmed
  carrier: uploaded_file
  format: image/png
  operation: parse_metadata
  attacker_influence: arbitrary_bytes
  crosses_process_boundary: true
  source: vendor
  evidence: https://vendor.example/advisory
  reviewed_at: 2026-07-09T00:00:00Z
\end{verbatim}

At minimum the record needs affected product/version scope, one or more
trigger profiles, provenance, publication time, assertion status, review
status, and supersession information. Scanner vendors should carry the record
next to the finding, but the authoritative assertion should normally originate
upstream with the vendor, CNA, or publisher that knows the vulnerable code.

\subsection{Agentic enrichment is a bridge, not the source of truth}
Agentic analysis can read advisories and propose a profile once per CVE or
product/version family. It cannot reconstruct facts that were never published.
The operational tradeoff is unavoidable:

\begin{itemize}[leftmargin=1.5em]
\item A conservative agent abstains frequently, leaving the scalability gap.
\item An aggressive agent increases coverage but creates unsafe false-negative
  risk when it concludes that no indirect path exists.
\item Human review improves reliability but recreates the manual burden if it
  is applied to every finding rather than selected CVE groups.
\end{itemize}

Agent output should therefore cite the exact source for every asserted field,
retain unknowns, record model and prompt versions, and require stronger review
for negative conclusions than for promotions. A proposed indirect trigger can
raise a CVE into the candidate set. An agent's unsupported conclusion that no
indirect trigger exists must not be treated as authoritative evidence.

The enrichment queue should be risk-ranked: KEVs, active exploitation, high
PAIN, high prevalence, and unresolved CVEs on components that commonly parse
untrusted content go first. Results are cached and invalidated when vendor or
threat intelligence changes.

\section{Optional Deployment-Evidence Refinement}
\label{sec:evidence-plane}
The current $C_0(a)$ predicate may come from a service archetype or governed
assertion; that is sufficient for Equation~\ref{eq:operational}. A provider or
tool developer may elect to refine $C_0(a)$ with a graph whose directed edges
describe actual content handoffs. Where used, declared policy and observed
telemetry have different roles:

\begin{itemize}[leftmargin=1.5em]
\item Infrastructure policy, NetworkPolicy, security groups, and service
  configuration describe which paths are permitted.
\item Flow, mesh, queue, trace, and host telemetry prove which paths have been
  observed and help discover undeclared edges.
\item Absence of observed traffic does not prove a future path is impossible.
  Only enforcement or a governed completeness assertion can rule it out.
\end{itemize}

An optional higher-resolution edge assertion can bind behavior and evidence to
the deployed artifact:

\begin{verbatim}
source: public-api
destination: order-database
channel: postgres
carries_internet_derived_data: true
transformations:
  - operation: parameterized_query
    effect: sql_structural_influence_to_data_only
evidence:
  source_commit: abc123
  artifact_digest: sha256:...
  sast_report: urn:...
  unit_test_report: urn:...
  integration_test_report: urn:...
  attested_by: ci.example.com
\end{verbatim}

The binding is the important part:

\[
\text{edge}\rightarrow\text{transformation assertion}
\rightarrow\text{test evidence}\rightarrow\text{source revision}
\rightarrow\text{deployed artifact}.
\]

Vulnerability-management platforms generally cannot consume this chain as a
portable finding-level reachability assertion today. The following are
ecosystem extension points, not provider maturity levels or requirements
imposed by this method:

\begin{center}
\renewcommand{\arraystretch}{1.08}
\footnotesize
\begin{tabular}{p{0.23\linewidth}p{0.66\linewidth}}
\toprule
\textbf{Profile} & \textbf{Capability} \\
\midrule
Recommended now & Scanner findings, CVSS, direct exposure, CVE-level
indirect-trigger promotion,
coarse content-receipt tags, and standing assessment/ConMon prevention
assertions. \\
Upstream reuse & Reusable CVE trigger profiles from authoritative or reviewed
enrichment. \\
Optional graph & Imported or dynamically derived application data-flow graphs. \\
Optional attestation & CI/CD-generated edge assertions and evidence bound to
deployed artifacts. \\
Optional automation & Continuous recomputation when vulnerabilities, code,
topology, controls, or
threat intelligence change. \\
\bottomrule
\end{tabular}
\end{center}

\section{Standing Prevention Evidence and Optional Transformations}
\label{sec:sanitization}
The operational $P_0(v,a)$ predicate does not depend on the optional
full-information graph. FedRAMP's Rev5 ConMon Playbook requires monthly scans
of web applications and all web interfaces or an approved sample, individual
tracking of findings, and assessor-validated scanner configuration. Providers
maintain the evidence and remediate findings; assessors verify the integrity
and results of provider-run scans. The 2026 VDR rules more generally require
persistent detection, validation, mitigation, and remediation.
\cite{fedramp-vdr,fedramp-conmon}

Together, those existing activities create a standing behavioral assertion for
the application and input-control classes actually covered by the validated
test program. $P_0$ makes that assertion explicit and revocable. It does not
claim that a clean DAST result discovers every trigger, proves every internal
path, or covers a vulnerability mechanism that the test plan does not exercise.
OWASP likewise cautions that automated scanners have coverage and false-negative
limits and should be used as part of a balanced testing program.
\cite{owasp-wstg}

Where a provider elects to support a higher-resolution transformation-based
NIRV assertion, the content algebra makes the same reasoning more precise. A
transformation does not remove a CWE label; it changes a capability carried by
the content.

\begin{itemize}[leftmargin=1.5em]
\item \textbf{Queue or storage handoff:} changes the carrier while normally
  preserving attacker influence and format.
\item \textbf{Parameterized SQL:} changes SQL influence from structural/control
  to data-only. The attacker string still reaches the database, but not as SQL
  grammar.
\item \textbf{Contextual output encoding:} changes HTML structural influence to
  text for the specific rendering context.
\item \textbf{Protocol termination:} changes protocol or version, such as
  \code{h2} to \code{http/1.1}.
\item \textbf{Schema validation:} restricts content to the language actually
  enforced by the schema; it does not generically remove deserialization or
  parser risk.
\item \textbf{Content disarm and reconstruction:} replaces unsafe document
  structures with a reconstructed format under an explicit policy.
\item \textbf{Unknown transformation:} preserves a confirmed applicable trigger
  capability until evidence establishes otherwise.
\end{itemize}

A content-trigger finding may earn NIRV operationally when the applicable
$P_0$ assertion is current. A higher-resolution claim may show that every path
carrying internet-derived input to the vulnerable operation contains a
transformation that prevents the specific trigger. The evidence hierarchy is:

\begin{center}
\renewcommand{\arraystretch}{1.2}
\small
\begin{tabular}{p{0.19\linewidth}p{0.68\linewidth}}
\toprule
\textbf{Weight} & \textbf{Examples} \\
\midrule
Operational & Assessor-validated scanner configuration and scope;
recurring DAST of the relevant interface and control class; relevant findings
remediated and retested; and no known bypass or material invalidating change. \\
Higher-resolution & Enforced path absence; protocol termination; affected code absent;
complete code/data-flow proof; typed or parameterized APIs; and targeted
regression or exploit tests bound to the deployed artifact. \\
Supporting & Scoped SAST, runtime observations, flow telemetry, unit and
integration tests, and general control assessment. \\
Insufficient alone & A clean scan with unknown scope or configuration, quiet
telemetry, an undocumented assertion, or absence of observed execution during
a finite window. \\
\bottomrule
\end{tabular}
\end{center}

FedRAMP's guidance expressly recognizes prevention before vulnerable
processing as relevant to IRV.\cite{fedramp-ver} The provider produces and
maintains the operational evidence and remediates findings. During assessment,
the 3PAO collects or reviews that evidence, validates its scope and method, and
verifies provider-run scan results. The assessor is not being asked to invent a
complete data-flow graph or perform routine CVE trigger enrichment on the
provider's behalf.

\begin{plain}
Regular DAST and remediation already say, in operational terms, ``the tested
input controls are working.'' The model uses that as a current, scoped
presumption, not as proof that every possible payload is safe. If a new CVE has
a trigger the testing did not cover, the presumption does not apply. The
optional graph can support the stronger statement that a named control blocks a
named trigger on every path, but the current method does not require providers
to build that graph.
\end{plain}

\section{Triage, Grouping, and Reporting}
\label{sec:triage}
The model retains more internal state than FedRAMP's binary reported field:

\begin{description}[leftmargin=10em,style=nextline]
\item[Confirmed IRV] A direct match exists, or a promoted indirect trigger
  applies and is not covered by current prevention evidence.
\item[Provisional IRV] An applicable trigger is known, but deployment or
  prevention evidence is incomplete.
\item[Evidenced NIRV] The direct branch is empty and each promoted indirect
  trigger is ruled out or covered by a current $P_0$ assertion.
\item[Indirect unknown] The upstream record does not establish whether an
  indirect trigger exists; the baseline has not promoted it.
\end{description}

The last state exposes the principal operational presumption. Defaulting every
unreviewed \code{AV:L/A/P} vulnerability to IRV would recover maximal
propagation and defeat scalability. The present profile therefore treats an
unknown indirect trigger as not promoted while preserving that status for
risk-ranked enrichment. This can temporarily miss a rare undocumented
indirect-trigger case. Threat-intelligence monitoring, KEV handling, vendor
updates, agentic proposals, analyst promotion, and re-evaluation are the
controls for that residual risk. The paper does not characterize the
presumption as proof.

The second operational presumption is that a current $P_0$ remains valid for
its documented interface, application archetype, and tested control class.
Recurring scans and remediation refresh it. A relevant finding, material
change, scope gap, new trigger technique, or contrary threat intelligence
revokes or narrows it. This is what makes convergence conditional and
defensible rather than a blanket claim that downstream content is sanitized.

Evaluation should be amortized across

\[
(\mathrm{CVE},\ \mathrm{product/version\ family},\
\mathrm{deployment\ archetype},\ \mathrm{input/control\ archetype})
\]

rather than repeated for every instance. One reviewed trigger profile can
serve thousands of scanner findings; one deployment archetype can supply the
same path and standing-control conclusion to a fleet of identically managed
assets. Instance-level exceptions remain possible and must break out of the
group.

The resulting IRV value plugs into the existing PAIN/LEV matrix without
changing PAIN:

\[
\text{column}=
\begin{cases}
\text{LEV+IRV} & \text{LEV}\wedge\mathrm{IRV}(v,a),\\
\text{LEV+NIRV} & \text{LEV}\wedge\neg\mathrm{IRV}(v,a),\\
\text{NLEV} & \neg\text{LEV}.
\end{cases}
\]

The provider may retain richer internal statuses and evidence while reporting
the binary value required by \code{VER-RPT-VDT}. Qualifying IAP/VPN edge
authorization is already reflected in $N(a)$. Source allowlisting, application
authentication, exploit maturity, and attacker prerequisites inform LEV or
attested mitigation rather than being forced into the reachability predicate.
Application authentication affects PAIN only when its enforced scope
demonstrably reduces the customer effect of successful exploitation.

\section{Worked Examples}
\label{sec:examples}
\begin{enumerate}[leftmargin=1.8em]
\item \textbf{Local vulnerability on a public web host.}
The CVE is \code{AV:L}, so it has no direct network descriptor and
$D(v,a)=0$. Its indirect status is retained separately. If $J(v)$ is unknown,
the current profile does not promote it; if later analysis shows that an
uploaded document invokes the vulnerable local parser, $J(v)$ becomes
confirmed and the indirect branch is evaluated across all affected assets. A
generic web scan does not set $P_0=1$ for that parser unless its validated scope
actually exercises the relevant upload and parser trigger class.

\item \textbf{SSH CVE on a public web host.}
The host exposes HTTPS publicly, but SSH is reachable only through a path that
the provider has evaluated outside the qualifying public population. The
finding's required surface contains SSH; $N(a)$ for the qualifying internet
path does not. Their intersection is empty even though the host is public.
Source attribution alone is not the proof; the recorded path interpretation
and enforcement are.

\item \textbf{HTTP/2 flaw behind protocol termination.}
A public load balancer accepts HTTP/2 but opens HTTP/1.1 connections to the
backend. The backend finding requires \code{h2}. With the load balancer as the
only path, the transformation removes \code{h2} from $N(\mathrm{backend})$ and
the finding is NIRV. The same flaw on a self-managed load balancer remains IRV.

\item \textbf{Administrative service behind a qualifying IAP.}
The IAP is the service's only ingress path and authorizes organizational
personnel before forwarding any request. The IAP remains internet-accessible;
the backend is removed from $N(a)$ and is NIRV for the direct branch. A login
implemented by the backend application would not produce the same result,
because the request reaches that component before the login logic runs.

\item \textbf{Known indirect logging trigger.}
Threat intelligence establishes that a crafted string can trigger a logging
library after crossing process boundaries. The CVE-level $J(v)$ is promoted
once. Every asset archetype tagged $C_0(a)=1$ and running the affected library
enters the IRV candidate group. Ordinary SQL-injection or cross-site-scripting
DAST coverage does not establish $P_0$ for a logging-expression trigger, so the
finding remains IRV unless a targeted control and its current evidence cover
the trigger. This is the current-profile treatment of the Log4Shell shape.

\item \textbf{SQL construction behind an application tier.}
Internet-derived values reach the database, but every application path uses
parameterized queries. The values therefore reach SQL only as data; they cannot
become SQL instructions. If the validated assessment and recurring DAST program
covers SQL injection on the relevant interface, relevant findings have been
remediated and retested, and no bypass or material change is known, then
$P_0(v,a)=1$ and the database finding is NIRV. The operational method does not
require a new proof of every query edge. A new SQL-injection finding, an
uncovered interface, or an unparameterized alternate path resets $P_0$ to zero.

\item \textbf{Installed but unobserved package.}
Runtime telemetry shows that a package did not execute during the observation
window. That is supporting evidence, not proof that its code is absent or can
never execute. CycloneDX \code{code\_not\_present} means the vulnerable code
has actually been removed or omitted; \code{code\_not\_reachable} requires a
durable execution-path conclusion, not merely a quiet window.\cite{cyclonedx}
\end{enumerate}

\section{Perimeter Prevention and Mitigation}
\label{sec:waf}
FedRAMP permits a finding to cease being considered IRV when a control
verifiably prevents the triggering payload before the vulnerable resource
processes it.\cite{fedramp-ver} A WAF virtual patch, API-gateway schema, mail
filter, or content-disarm boundary can qualify when it covers the actual
trigger and every relevant path.

The stricter operational posture is still to preserve intrinsic reachability
and record prevention separately. WAF rules are tuned, bypassed, moved into
detection-only mode, and sometimes removed during incidents. Keeping the
finding visible as intrinsically reachable while attaching a signed VEX or
other mitigation assertion produces a record that degrades safely when the
control fails. CycloneDX provides machine-readable impact-analysis states and
justifications for such assertions.\cite{cyclonedx}

This does not conflict with $P_0$, which is refreshed through validated
assessment and ConMon for scoped application and input-control classes. The
paper keeps mutable, CVE-specific perimeter rules in the mitigation record
unless the provider establishes the same scope, currency, and invalidation bar.

Known Exploited Vulnerabilities retain the remediation expectation in
\code{VDR-TFR-KEV} even when fully mitigated.\cite{fedramp-vdr} The provider
should therefore track reachability, prevention, mitigation, disposition, and
the KEV clock as separate facts.

\section{Limitations and Validation Requirements}
\label{sec:limitations}
This is a triage architecture, not a proof that all transitive triggers are
known. Its principal limitation is deliberate: an indirect trigger whose
semantics are absent from upstream metadata and not promoted by enrichment can
remain outside $R_{\mathrm{operational}}$. The method trades exhaustive recall
for operational scale and makes that trade visible through the \code{unknown}
state. Its second limitation is that $P_0$ is only as reliable as the validated
scope, currency, coverage, remediation, and change-invalidation discipline of
the assessment and ConMon program.

The claim that undocumented indirect-trigger vulnerabilities are exceptional
should be measured rather than asserted. A reference implementation should be
tested on a reviewed corpus containing direct protocol flaws, connection-bound
flaws, local privilege vulnerabilities, file and media parsers, deserializers,
logging and expression systems, and known transitive cases. At minimum it
should report:

\begin{itemize}[leftmargin=1.5em,nosep]
\item false-negative and false-positive rates for indirect-trigger promotion;
\item the proportion of CVEs left unknown;
\item agreement across analysts, agents, model versions, and repeated runs;
\item provenance coverage and the authority of cited sources;
\item enrichment cost and latency per previously unseen CVE; and
\item time from new threat intelligence to fleet-wide reclassification;
\item DAST interface and trigger-class coverage, including authenticated paths;
  and
\item time from a relevant finding or material change to invalidation of $P_0$.
\end{itemize}

When a provider elects to use the optional evidence graph to earn a narrower
finding-level exclusion, that claim needs validation. Edge completeness,
binding to deployed artifacts, change invalidation, and alternate-path
discovery are security properties. A signed assertion is useful only if its
scope is exact and its inputs remain true. None of this converts the optional
reference form into a baseline requirement for providers using
Equation~\ref{eq:operational}.

\section{Conclusion}
FedRAMP's finding-level internet-reachability requirement exposes a metadata
problem upstream of any individual provider. Direct reachability can be
computed reasonably well from network, runtime, and scanner evidence. The
primary obstacle to scalable transitive classification is that today's
vulnerability sources and scanners do not reliably identify which CVEs have an
indirect internet-content trigger. More detailed deployment evidence can refine
where a known trigger applies, but a complete application-wide evidence graph
is not required by the operational method recommended here.

The operational profile in this paper is intentionally modest. It computes
the direct branch, uses CVSS attack vector only for that branch, promotes known
indirect triggers once per CVE, and joins them to coarse content-receipt
archetypes. It also reuses assessment- and ConMon-backed prevention assertions
for the input-control classes already tested. The resulting operational set
therefore converges toward the directly accessible surface, plus known indirect
triggers whose path or prevention evidence is missing, stale, failed, or out of
scope. That is a conditional operational convergence, not a semantic claim that
internet reachability and internet accessibility are identical.

Two explicit presumptions make the method scalable. Unknown indirect triggers
are not promoted until upstream intelligence or review identifies them, and a
current $P_0$ assertion is trusted within its validated scope until contrary
evidence or change revokes it. The first creates residual risk from missing
trigger metadata. The second creates residual risk from incomplete testing or
stale control evidence. Threat intelligence, KEV handling, recurring DAST,
finding remediation and retest, change detection, agentic proposals, selective
analyst review, and rapid reclassification control those risks.

Equation~\ref{eq:operational} is the paper's recommendation for that current
environment. The full-information form is non-normative: it documents an
interoperability option if vendors, CNAs, scanners, or provider-selected tools
later supply richer inputs. It does not establish that FedRAMP expects
providers to build a complete data-flow attestation system now, and its
existence does not make the present classifier deficient for its declared
triage purpose. Existing DAST and remediation evidence is sufficient for the
paper's scoped standing-control presumption; the optional graph would improve
precision, not create the premise from nothing. The durable near-term
improvement is upstream: publish
provenance-backed trigger profiles once and carry them through scanners so
providers can promote exceptional indirect cases consistently. Agentic
enrichment can bridge that gap, but it does not become the authority for facts
no source has disclosed.

\begin{thebibliography}{9}
\bibitem{fedramp-definitions}
FedRAMP, \emph{FedRAMP Definitions: Internet-Reachable Vulnerability},
Consolidated Rules for 2026.\\
\href{https://www.fedramp.gov/2026/definitions/}{FedRAMP definitions}

\bibitem{fedramp-ver}
FedRAMP, \emph{Vulnerability Evaluation and Reporting}, Consolidated Rules
for 2026, including \code{VER-EVA-EIR}, \code{VER-EVA-GRV}, and
\code{VER-RPT-VDT}.\\
\href{https://www.fedramp.gov/2026/reference/rev5/b/vulnerability-evaluation-and-reporting/}{FedRAMP VER reference}

\bibitem{fedramp-vdr}
FedRAMP, \emph{Vulnerability Detection and Response}, Consolidated Rules for
2026, including \code{VDR-TFR-PVR} and \code{VDR-TFR-KEV}.\\
\href{https://www.fedramp.gov/2026/reference/rev5/d/vulnerability-detection-and-response/}{FedRAMP VDR reference}

\bibitem{fedramp-conmon}
FedRAMP, \emph{Continuous Monitoring Playbook}, Version 1.0, November 17,
2025, vulnerability scanning and assessor responsibilities.\\
\href{https://www.fedramp.gov/resources/documents/Continuous_Monitoring_Playbook.pdf}{FedRAMP ConMon Playbook}

\bibitem{cvss-user-guide}
Forum of Incident Response and Security Teams (FIRST), \emph{CVSS v4.0 User
Guide}, ``Attack Vector Considerations.''\\
\href{https://www.first.org/cvss/user-guide.html}{FIRST CVSS v4.0 User Guide}

\bibitem{cyclonedx}
OWASP CycloneDX, \emph{CycloneDX v1.7 Specification Reference}, impact
analysis states and justifications.\\
\href{https://cyclonedx.org/docs/1.7/}{CycloneDX v1.7 reference}

\bibitem{owasp-wstg}
OWASP Foundation, \emph{Web Security Testing Guide v4.2}, introduction and
limitations of automated testing.\\
\href{https://owasp.org/www-project-web-security-testing-guide/v42/2-Introduction/README}{OWASP WSTG v4.2 introduction}
\end{thebibliography}

\section*{Glossary of Abbreviations}
\begin{center}
\footnotesize
\begin{tabular}{>{\ttfamily}p{0.08\linewidth}p{0.34\linewidth}>{\ttfamily}p{0.08\linewidth}p{0.34\linewidth}}
3PAO & Third-Party Assessment Organization & ALPN & Application-Layer Protocol Negotiation \\
CISA & Cybersecurity and Infrastructure Security Agency & CNA & CVE Numbering Authority \\
ConMon & Continuous Monitoring & CPE & Common Platform Enumeration \\
CVE & Common Vulnerabilities and Exposures & CVSS & Common Vulnerability Scoring System \\
CWE & Common Weakness Enumeration & DAST & Dynamic Application Security Testing \\
IRV & Internet-Reachable Vulnerability & KEV & Known Exploited Vulnerability \\
LEV & Likely Exploitable Vulnerability & NIRV & Not Internet-reachable Vulnerability \\
PAIN & Potential Agency Impact N-rating & SAST & Static Application Security Testing \\
VDR & Vulnerability Detection and Response & VER & Vulnerability Evaluation and Reporting \\
\end{tabular}
\end{center}

\end{document}
